\documentclass[11pt]{article}
\usepackage[utf8]{inputenc}
\usepackage{amsmath, amssymb}
\usepackage{geometry}
\usepackage{listings}
\usepackage{xcolor}
\usepackage{tabularx}
\usepackage{booktabs}
\usepackage{hyperref}

\geometry{a4paper, margin=1in}

\definecolor{codegray}{rgb}{0.5,0.5,0.5}
\definecolor{codepurple}{rgb}{0.58,0,0.82}
\definecolor{backcolour}{rgb}{0.95,0.95,0.92}

\lstset{
    language=Python,
    backgroundcolor=\color{backcolour},
    commentstyle=\color{codegray},
    keywordstyle=\color{magenta},
    stringstyle=\color{codepurple},
    basicstyle=\ttfamily\small,
    breaklines=true,
    showstringspaces=false
}

\title{Physics Lab: Automatic Differentiation \\ \large Problem Set \& Implementation Guide}
\author{VB}
\date{\today}

\begin{document}

\maketitle

\section*{Problem 0: Start the engine}

Consider the files given.
\begin{center}
\renewcommand{\arraystretch}{1.5}
\begin{tabularx}{\textwidth}{@{} >{\bfseries\raggedright\arraybackslash}p{5.5cm} X @{}}
\toprule
\textcolor{blue}{Filename} & \textcolor{blue}{Purpose} \\ \midrule
\multicolumn{2}{@{}l}{\textbf{Direct Interaction}} \\ \midrule
\lstinline|implementation_tasks.py| & \textcolor{red}{Here you write your code.} Contains the core math logic for Dual Numbers ($a + b\epsilon$) intended for student implementation. \\
\lstinline|lab_dashboard.py| & \textcolor{red}{This file you run.} A dedicated educational sandbox for debugging individual math levels through an interactive UI. \\
\lstinline|main_demo.py| & \textbf{The AutoDiff Engine.} The main standalone demo where you can write multiline Python code with cycles and conditionals, and plot their exact derivatives. \\ \midrule
\multicolumn{2}{@{}l}{\textbf{Under the Hood}} \\ \midrule
\lstinline|levels/base_level.py| & Base class for graphical debuggers.\\
\lstinline|levels/level_*.py| & Subproblem interfaces for arithmetic, Taylor truncation, and transcendental functions. \\
\bottomrule
\end{tabularx}
\end{center}

\textbf{Prerequisites:} You will need to install the dependencies for this lab. Run \lstinline|pip install -r requirements.txt| before starting.

Now use any editor to open the file \lstinline|implementation_tasks.py|. You should see API for functions you should implement. You can add your own \lstinline|unittest.TestCase| classes at the bottom of \lstinline|implementation_tasks.py| and run \lstinline|python implementation_tasks.py| to execute them.

\subsubsection*{Self-Testing (Optional)}
At the bottom of \lstinline|implementation_tasks.py|, you will find a block that looks like this:
\begin{lstlisting}[language=Python]
if __name__ == '__main__':
    import unittest
    # Add your own unittest.TestCase classes here...
    unittest.main(verbosity=2)
\end{lstlisting}
This allows you to write your own automated unit tests to verify your code before running the visual dashboard. If you'd like to test your functions with specific inputs, simply define a class inheriting from \lstinline|unittest.TestCase| above the \lstinline|unittest.main()| call. For example, if you wanted to test simple math:
\begin{lstlisting}[language=Python]
class TestMyMath(unittest.TestCase):
    def test_addition(self):
        result = 2 + 2
        self.assertEqual(result, 4, "Addition is broken!")
\end{lstlisting}
You can run these tests anytime by executing \lstinline|python implementation_tasks.py| in your terminal. This is completely optional, but highly recommended for debugging tricky edge cases.
\begin{lstlisting}
import numpy as np
import math
import unittest

class Dual:
    """Dual number: a + b*eps where eps^2 = 0."""
    def __init__(self, real, dual):
        self.real = float(real)
        self.dual = float(dual)

    def __add__(self, other):
        pass

<<<< MORE CODE >>>>
\end{lstlisting}

\section*{Problem 1: Dual Number Arithmetic}

\subsubsection*{Mathematical part}
\textbf{Given:} Two dual numbers $X = a + b\epsilon$ and $Y = c + d\epsilon$, where $\epsilon^2 = 0$. \\
\textbf{Derive:} The formulas for addition, subtraction, multiplication, and division.
For example, for multiplication:
$$ X \cdot Y = (a + b\epsilon)(c + d\epsilon) = ac + (ad + bc)\epsilon + bd\epsilon^2 $$
Since $\epsilon^2 = 0$, this simplifies to:
$$ X \cdot Y = ac + (ad + bc)\epsilon $$

\subsubsection*{Implementation part}
\textbf{Implement:} The standard Python magic methods \lstinline|__add__|, \lstinline|__sub__|, \lstinline|__mul__|, \lstinline|__truediv__|, and \lstinline|__pow__| in the \lstinline|Dual| class. Make sure to handle operations where \lstinline|other| is a standard float/int.

\textbf{Test:} Run \lstinline|lab_dashboard.py| and use "L1: Elementary Arithmetic" to check your symbolic derivations. Ensure unittests pass.

\section*{Problem 2: Transcendental Functions (Chain Rule)}

\subsubsection*{Mathematical part}
\textbf{Given:} A dual number $x = a + b\epsilon$. \\
\textbf{Derive:} The Taylor expansion of $f(x)$ centered around $a$. Show that because $\epsilon^2 = 0$, the infinite series truncates exactly to $f(a) + f'(a)b\epsilon$. Use this to derive the definitions of $\sin(x)$, $\cos(x)$, $\tan(x)$, $\exp(x)$, and $\log(x)$ for Dual numbers.

\subsubsection*{Implementation part}
\textbf{Implement:} The standalone functions \lstinline|sin(x)|, \lstinline|cos(x)|, \lstinline|tan(x)|, \lstinline|exp(x)|, and \lstinline|log(x)|. They should check if \lstinline|x| is a \lstinline|Dual| number; if so, apply your derived formulas. If not, fallback to standard \lstinline|math| functions.

\textbf{Test:} Use "L3: Elementary Functions" in the dashboard. Ensure that the dual component ($\epsilon$) plotted perfectly matches the analytical derivative.

\section*{Problem 3: The AutoDiff Demo}
Once your core math engine is passing all tests, launch the Main Demo using the dashboard or by running \lstinline|python3 main_demo.py|. Write an arbitrary, complex iterative algorithm using standard Python logic (loops, conditionals) in the text box. The underlying Dual number engine will flow through your logic and plot the perfect analytical derivative of your algorithm automatically!

\section*{\textcolor{red}{Submission}}
Submit the file \texttt{implementation\_tasks.py}.

\end{document}
